%! Setting Up a Test for a Population Proportion — Unit 6, Lesson 6.4 — Warm-Up
\documentclass[10pt]{article}
\usepackage{apstats-article}
\usepackage{apstats-boxes}
\newcommand{\TallMath}[1]{$\displaystyle #1\rule[-1.4em]{0pt}{3.2em}$}

\begin{document}

\pageheader{Unit 6, Lesson 6.4}{Warm-Up}
\namedateperiod

\vspace{0.15in}

A 95\% CI for the proportion of students who own a pet is $(0.43, 0.57)$.

\begin{enumerate}
  \item Write a complete interpretation of this CI.
        \writelines{2}

  \item A school board member claims 60\% of students own a pet. Is this plausible
        based on the interval? Explain.
        \writelines{2}

  \item Suppose instead of estimating $p$, a researcher wants to \emph{test} whether the
        true proportion differs from 0.50. Write the null and alternative hypotheses.

        $H_0$: \blank{5cm}  \quad $H_a$: \blank{5cm}

  \item In the CI, you used $\hat p$ in the Large Counts condition. In a hypothesis test,
        which value would you use instead, and why?
        \writelines{2}
\end{enumerate}

\end{document}
